% ========================================
% 講義MODE プリアンブル（A4 縦 × 最大4ページ → B4 横・表裏1枚に左右2面で縮小印刷）
%   1・2ページ目 → 表面，3・4ページ目 → 裏面
% 使い方: \documentclass[14pt]{jsarticle}
%         % ========================================
% 講義MODE プリアンブル（A4 縦 × 最大4ページ → B4 横・表裏1枚に左右2面で縮小印刷）
%   1・2ページ目 → 表面，3・4ページ目 → 裏面
% 使い方: \documentclass[14pt]{jsarticle}
%         % ========================================
% 講義MODE プリアンブル（A4 縦 × 最大4ページ → B4 横・表裏1枚に左右2面で縮小印刷）
%   1・2ページ目 → 表面，3・4ページ目 → 裏面
% 使い方: \documentclass[14pt]{jsarticle}
%         % ========================================
% 講義MODE プリアンブル（A4 縦 × 最大4ページ → B4 横・表裏1枚に左右2面で縮小印刷）
%   1・2ページ目 → 表面，3・4ページ目 → 裏面
% 使い方: \documentclass[14pt]{jsarticle}
%         \input{lecture-preamble.tex}
% ※ B4 に縮小（約87%）すると，14pt は実寸で約12pt になる
% ※ jsarticle は \mag で拡大するので，紙面の寸法は truemm で指定する．
%    TikZ の cm も拡大される（14pt では 1cm → 実寸 1.4cm）点に注意
% ========================================
% a4paper は \mag で拡大されてしまうため，用紙も truemm で実寸指定する
\usepackage[dvipdfm, paperwidth=210truemm, paperheight=297truemm,
            margin=15truemm]{geometry}

\input{common-preamble.tex}

% A4 2ページに収めるため，共通設定（1.25）より行間を少し詰める
\linespread{1.15}
\pagestyle{empty}

% ========================================
% ヘッダ（1ページ目の先頭）
% 使い方: \lectureheader{科目}{章番号}{章の名前}{節番号}{節の名前}{項番号-小単元}{小単元の名前}
%           {教科書のページ}{本時の目標}
%   例: \lectureheader{数学Ⅱ}{5}{指数関数と対数関数}{1}{指数関数}{2-A}{指数関数 $y=a^x$ のグラフ}
%         {150--151}{指数関数のグラフを描けるようになろう}
%   項番号-小単元は「2-A」のほか「2」だけでもよい（カウンターには先頭の数字を使う）
%   教科書のページを空 {} にすると，手書き用の空欄になる
%   タイトルと教科書のページが1行に収まらないときは，タイトルを横に縮めて収める
% ========================================
\newsavebox{\lectitlebox}
\newsavebox{\lecpagebox}
\newlength{\lectitlewidth}
\def\lecitemnumber#1-#2\relax{#1}
\newcommand{\lectureheader}[9]{%
  \setcounter{mathchapter}{#2}%
  \setcounter{section}{#4}%
  \setcounter{subsection}{\lecitemnumber#6-\relax}%
  \noindent{\bfseries #1}\quad 第#2章\quad #3\hfill
  \rule{0pt}{1.4em}%
  \underline{\hspace{1.5zw}}年\underline{\hspace{1.5zw}}組\underline{\hspace{2zw}}番\quad
  氏名\underline{\hspace{8zw}}\par
  \vspace{0.3em}
  \sbox{\lectitlebox}{\boldmath{\Large\bfseries 第#4節\enspace #5}\hspace{1zw}{\large\bfseries #6\enspace #7}}%
  \sbox{\lecpagebox}{教科書 p.\,\ifx\relax#8\relax\hspace*{2zw}\else#8\fi}%
  \setlength{\lectitlewidth}{\dimexpr\linewidth-\wd\lecpagebox-1zw\relax}%
  \noindent
  \ifdim\wd\lectitlebox>\lectitlewidth
    \resizebox{\lectitlewidth}{!}{\usebox{\lectitlebox}}%
  \else
    \usebox{\lectitlebox}%
  \fi
  \hfill\usebox{\lecpagebox}\par
  \vspace{-0.3em}
  \noindent\rule{\linewidth}{0.8pt}\par
  \vspace{0.3em}
  \begin{itembox}[l]{\textbf{本時の目標}}
    \centering\bfseries\boldmath #9
  \end{itembox}
  \vspace{0.3em}
}

% 次のページへ移る（1→2 は表面の右，2→3 は裏面の左，3→4 は裏面の右）
\newcommand{\nextpage}{\newpage}

% 小単元（教科書の「A」「B」…）: \topic{A}{平均変化率}
\newcommand{\topic}[2]{%
  \par\vspace{0.3em}\noindent{\Large\bfseries #1\quad #2}\par\vspace{0.3em}%
}

% 見出し: \heading{見出し}
\newcommand{\heading}[1]{%
  \par\noindent{\large\bfseries #1}\par\vspace{0.3em}%
}

% 前時の復習（導入）: \begin{reviewbox}{中学数学：1次関数の変化の割合} ... \end{reviewbox}
\newenvironment{reviewbox}[1]{%
  \begin{itembox}[l]{\textbf{前時の復習}\quad #1}
}{\end{itembox}}

% 本時のまとめ
\newenvironment{summarybox}{%
  \begin{itembox}[l]{\textbf{本時のまとめ}}
}{\end{itembox}}

% 左右配置（左：説明・計算，右：図）
% 使い方: \sidebyside[左の幅の割合]{左}{右}
\newcommand{\sidebyside}[3][0.5]{%
  \par\noindent
  \begin{minipage}[t]{#1\linewidth}\vspace{0pt}#2\end{minipage}\hfill
  \begin{minipage}[t]{\dimexpr0.98\linewidth-#1\linewidth\relax}%
    \vspace{0pt}\centering #3\end{minipage}\par
}

% 次回予告（最終ページの下端）
\newcommand{\nexttime}[1]{%
  \par\noindent\parbox{\linewidth}{%
    \rule{\linewidth}{0.4pt}\par
    次回：#1}\par
}

% ※ B4 に縮小（約87%）すると，14pt は実寸で約12pt になる
% ※ jsarticle は \mag で拡大するので，紙面の寸法は truemm で指定する．
%    TikZ の cm も拡大される（14pt では 1cm → 実寸 1.4cm）点に注意
% ========================================
% a4paper は \mag で拡大されてしまうため，用紙も truemm で実寸指定する
\usepackage[dvipdfm, paperwidth=210truemm, paperheight=297truemm,
            margin=15truemm]{geometry}

% ========================================
% 共通プリアンブル（講義MODE・演習MODE）
% パッケージ／解答切り替え／番号／空欄／定義・公式・定理の囲み
% ========================================
\usepackage{amsmath, amssymb, mathtools, bm}
\usepackage{ascmac}
\usepackage[dvipdfmx]{graphicx}
\usepackage[dvipdfmx]{xcolor}
\usepackage{tikz}
\usetikzlibrary{calc, arrows.meta}
\usepackage[most]{tcolorbox}
\usepackage{enumitem}
\usepackage{tasks}

\linespread{1.25}
\everymath{\displaystyle}
\setlength{\parindent}{0pt}
\setlist[enumerate,1]{label=（\arabic*）, leftmargin=*, labelsep=0.5em}

% ========================================
% 解答表示の切り替え（既定は教師用．make student で \studentmode が定義され生徒用になる）
% ========================================
\newif\ifshowanswer
\showanswertrue
\ifdefined\studentmode\showanswerfalse\fi

% ========================================
% 番号（「章.節.番号」形式，節ごとにリセット）
% ========================================
% jsarticle には章がないので，章番号は専用カウンタで持つ
\newcounter{mathchapter}
\newcounter{definition}[section]
\newcounter{formula}[section]
\newcounter{theorem}[section]
\renewcommand{\thedefinition}{\arabic{mathchapter}.\arabic{section}.\arabic{definition}}
\renewcommand{\theformula}{\arabic{mathchapter}.\arabic{section}.\arabic{formula}}
\renewcommand{\thetheorem}{\arabic{mathchapter}.\arabic{section}.\arabic{theorem}}

% ========================================
% 書き込み用コマンド
% ========================================
% 解答は青字・太字（数式も和文も太字になるよう \bfseries\boldmath を使う）
% 色は外側から \textcolor で掛けると数式中の \text{} の和文に効かないため，
% 箱の内側で \color を指定する
\newcommand{\answerstyle}[1]{{\color{blue}\bfseries\boldmath #1}}

% 空欄の左右それぞれの余白
\newcommand{\fitblankpad}{2zw}

% 空欄（答えの幅＋左右の余白で下線を引く）
% 使い方: \fitblankbf[]{答え}  ※文字列は \text{...} で囲む
\newsavebox{\fitblanksavebox}
\newcommand{\fitblankbf}[2][]{%
  \sbox{\fitblanksavebox}{\answerstyle{$\displaystyle #2$}}%
  \underline{\makebox[\dimexpr\wd\fitblanksavebox+\fitblankpad*2\relax][c]{%
    \ifshowanswer\usebox{\fitblanksavebox}%
    \else\phantom{\usebox{\fitblanksavebox}}\fi}}%
}

% 下線なしの空欄（分数の分子・分母など，下線と分数線が重なる場所で使う）
% 使い方: \frac{\fitblankbox{y \text{ の変化量}}}{\fitblankbox{x \text{ の変化量}}}
\newcommand{\fitblankbox}[1]{%
  \sbox{\fitblanksavebox}{\answerstyle{$\displaystyle #1$}}%
  \makebox[\dimexpr\wd\fitblanksavebox+\fitblankpad*2\relax][c]{%
    \ifshowanswer\usebox{\fitblanksavebox}%
    \else\phantom{\usebox{\fitblanksavebox}}\fi}%
}

% 数式中の計算欄（下線なし，答えの幅＋左右2em を確保）
% 使い方: $= \answermath{3 - 1 = 2}$
\newcommand{\answermath}[1]{%
  \ifshowanswer\mbox{\answerstyle{$\displaystyle #1$}}%
  \else\hspace{2em}\phantom{#1}\hspace{2em}\fi
}

% 記述・計算スペース（薄い枠で高さを確保）
% 使い方: \writespace[高さ]{解答例}
\newcommand{\writespace}[2][5em]{%
  \begin{tcolorbox}[enhanced, colback=white, colframe=black!35,
      boxrule=0.4pt, sharp corners, height=#1, valign=center,
      left=8pt, right=8pt]
    \ifshowanswer\answerstyle{#2}\fi
  \end{tcolorbox}%
}

% ========================================
% 定義（赤系）・公式（青系）・定理（青系・太枠）
% 使い方: \begin{definitionbox}{用語} ... \end{definitionbox}
% ========================================
\tcbset{
  mathbox/.style={
    enhanced, sharp corners,
    left=10pt, right=10pt, top=8pt, bottom=8pt,
    fonttitle=\bfseries,
  }
}
\newtcolorbox{definitionbox}[1]{
  mathbox, colback=red!5, colframe=red!70!black, boxrule=0.8pt,
  code={\refstepcounter{definition}},
  title={定義 \thedefinition\quad #1}
}
\newtcolorbox{formulabox}[1]{
  mathbox, colback=blue!5, colframe=blue!70!black, boxrule=0.8pt,
  code={\refstepcounter{formula}},
  title={公式 \theformula\quad #1}
}
\newtcolorbox{theorembox}[1]{
  mathbox, colback=blue!5, colframe=blue!70!black, boxrule=1.2pt,
  code={\refstepcounter{theorem}},
  title={定理 \thetheorem\quad #1}
}
% テクニック（緑系）
% 使い方: \begin{techniquebox}[ラベル]{見出し} ... \end{techniquebox}
\newcounter{technique}[section]
\renewcommand{\thetechnique}{\arabic{mathchapter}.\arabic{section}.\arabic{technique}}
\newtcolorbox{techniquebox}[2][]{
  mathbox, colback=green!5, colframe=green!70!black, boxrule=0.8pt,
  code={\refstepcounter{technique}\if\relax\detokenize{#1}\relax\else\label{#1}\fi},
  title={Technique \thetechnique\if\relax\detokenize{#2}\relax\else\quad\textbf{#2}\fi}
}


% A4 2ページに収めるため，共通設定（1.25）より行間を少し詰める
\linespread{1.15}
\pagestyle{empty}

% ========================================
% ヘッダ（1ページ目の先頭）
% 使い方: \lectureheader{科目}{章番号}{章の名前}{節番号}{節の名前}{項番号-小単元}{小単元の名前}
%           {教科書のページ}{本時の目標}
%   例: \lectureheader{数学Ⅱ}{5}{指数関数と対数関数}{1}{指数関数}{2-A}{指数関数 $y=a^x$ のグラフ}
%         {150--151}{指数関数のグラフを描けるようになろう}
%   項番号-小単元は「2-A」のほか「2」だけでもよい（カウンターには先頭の数字を使う）
%   教科書のページを空 {} にすると，手書き用の空欄になる
%   タイトルと教科書のページが1行に収まらないときは，タイトルを横に縮めて収める
% ========================================
\newsavebox{\lectitlebox}
\newsavebox{\lecpagebox}
\newlength{\lectitlewidth}
\def\lecitemnumber#1-#2\relax{#1}
\newcommand{\lectureheader}[9]{%
  \setcounter{mathchapter}{#2}%
  \setcounter{section}{#4}%
  \setcounter{subsection}{\lecitemnumber#6-\relax}%
  \noindent{\bfseries #1}\quad 第#2章\quad #3\hfill
  \rule{0pt}{1.4em}%
  \underline{\hspace{1.5zw}}年\underline{\hspace{1.5zw}}組\underline{\hspace{2zw}}番\quad
  氏名\underline{\hspace{8zw}}\par
  \vspace{0.3em}
  \sbox{\lectitlebox}{\boldmath{\Large\bfseries 第#4節\enspace #5}\hspace{1zw}{\large\bfseries #6\enspace #7}}%
  \sbox{\lecpagebox}{教科書 p.\,\ifx\relax#8\relax\hspace*{2zw}\else#8\fi}%
  \setlength{\lectitlewidth}{\dimexpr\linewidth-\wd\lecpagebox-1zw\relax}%
  \noindent
  \ifdim\wd\lectitlebox>\lectitlewidth
    \resizebox{\lectitlewidth}{!}{\usebox{\lectitlebox}}%
  \else
    \usebox{\lectitlebox}%
  \fi
  \hfill\usebox{\lecpagebox}\par
  \vspace{-0.3em}
  \noindent\rule{\linewidth}{0.8pt}\par
  \vspace{0.3em}
  \begin{itembox}[l]{\textbf{本時の目標}}
    \centering\bfseries\boldmath #9
  \end{itembox}
  \vspace{0.3em}
}

% 次のページへ移る（1→2 は表面の右，2→3 は裏面の左，3→4 は裏面の右）
\newcommand{\nextpage}{\newpage}

% 小単元（教科書の「A」「B」…）: \topic{A}{平均変化率}
\newcommand{\topic}[2]{%
  \par\vspace{0.3em}\noindent{\Large\bfseries #1\quad #2}\par\vspace{0.3em}%
}

% 見出し: \heading{見出し}
\newcommand{\heading}[1]{%
  \par\noindent{\large\bfseries #1}\par\vspace{0.3em}%
}

% 前時の復習（導入）: \begin{reviewbox}{中学数学：1次関数の変化の割合} ... \end{reviewbox}
\newenvironment{reviewbox}[1]{%
  \begin{itembox}[l]{\textbf{前時の復習}\quad #1}
}{\end{itembox}}

% 本時のまとめ
\newenvironment{summarybox}{%
  \begin{itembox}[l]{\textbf{本時のまとめ}}
}{\end{itembox}}

% 左右配置（左：説明・計算，右：図）
% 使い方: \sidebyside[左の幅の割合]{左}{右}
\newcommand{\sidebyside}[3][0.5]{%
  \par\noindent
  \begin{minipage}[t]{#1\linewidth}\vspace{0pt}#2\end{minipage}\hfill
  \begin{minipage}[t]{\dimexpr0.98\linewidth-#1\linewidth\relax}%
    \vspace{0pt}\centering #3\end{minipage}\par
}

% 次回予告（最終ページの下端）
\newcommand{\nexttime}[1]{%
  \par\noindent\parbox{\linewidth}{%
    \rule{\linewidth}{0.4pt}\par
    次回：#1}\par
}

% ※ B4 に縮小（約87%）すると，14pt は実寸で約12pt になる
% ※ jsarticle は \mag で拡大するので，紙面の寸法は truemm で指定する．
%    TikZ の cm も拡大される（14pt では 1cm → 実寸 1.4cm）点に注意
% ========================================
% a4paper は \mag で拡大されてしまうため，用紙も truemm で実寸指定する
\usepackage[dvipdfm, paperwidth=210truemm, paperheight=297truemm,
            margin=15truemm]{geometry}

% ========================================
% 共通プリアンブル（講義MODE・演習MODE）
% パッケージ／解答切り替え／番号／空欄／定義・公式・定理の囲み
% ========================================
\usepackage{amsmath, amssymb, mathtools, bm}
\usepackage{ascmac}
\usepackage[dvipdfmx]{graphicx}
\usepackage[dvipdfmx]{xcolor}
\usepackage{tikz}
\usetikzlibrary{calc, arrows.meta}
\usepackage[most]{tcolorbox}
\usepackage{enumitem}
\usepackage{tasks}

\linespread{1.25}
\everymath{\displaystyle}
\setlength{\parindent}{0pt}
\setlist[enumerate,1]{label=（\arabic*）, leftmargin=*, labelsep=0.5em}

% ========================================
% 解答表示の切り替え（既定は教師用．make student で \studentmode が定義され生徒用になる）
% ========================================
\newif\ifshowanswer
\showanswertrue
\ifdefined\studentmode\showanswerfalse\fi

% ========================================
% 番号（「章.節.番号」形式，節ごとにリセット）
% ========================================
% jsarticle には章がないので，章番号は専用カウンタで持つ
\newcounter{mathchapter}
\newcounter{definition}[section]
\newcounter{formula}[section]
\newcounter{theorem}[section]
\renewcommand{\thedefinition}{\arabic{mathchapter}.\arabic{section}.\arabic{definition}}
\renewcommand{\theformula}{\arabic{mathchapter}.\arabic{section}.\arabic{formula}}
\renewcommand{\thetheorem}{\arabic{mathchapter}.\arabic{section}.\arabic{theorem}}

% ========================================
% 書き込み用コマンド
% ========================================
% 解答は青字・太字（数式も和文も太字になるよう \bfseries\boldmath を使う）
% 色は外側から \textcolor で掛けると数式中の \text{} の和文に効かないため，
% 箱の内側で \color を指定する
\newcommand{\answerstyle}[1]{{\color{blue}\bfseries\boldmath #1}}

% 空欄の左右それぞれの余白
\newcommand{\fitblankpad}{2zw}

% 空欄（答えの幅＋左右の余白で下線を引く）
% 使い方: \fitblankbf[]{答え}  ※文字列は \text{...} で囲む
\newsavebox{\fitblanksavebox}
\newcommand{\fitblankbf}[2][]{%
  \sbox{\fitblanksavebox}{\answerstyle{$\displaystyle #2$}}%
  \underline{\makebox[\dimexpr\wd\fitblanksavebox+\fitblankpad*2\relax][c]{%
    \ifshowanswer\usebox{\fitblanksavebox}%
    \else\phantom{\usebox{\fitblanksavebox}}\fi}}%
}

% 下線なしの空欄（分数の分子・分母など，下線と分数線が重なる場所で使う）
% 使い方: \frac{\fitblankbox{y \text{ の変化量}}}{\fitblankbox{x \text{ の変化量}}}
\newcommand{\fitblankbox}[1]{%
  \sbox{\fitblanksavebox}{\answerstyle{$\displaystyle #1$}}%
  \makebox[\dimexpr\wd\fitblanksavebox+\fitblankpad*2\relax][c]{%
    \ifshowanswer\usebox{\fitblanksavebox}%
    \else\phantom{\usebox{\fitblanksavebox}}\fi}%
}

% 数式中の計算欄（下線なし，答えの幅＋左右2em を確保）
% 使い方: $= \answermath{3 - 1 = 2}$
\newcommand{\answermath}[1]{%
  \ifshowanswer\mbox{\answerstyle{$\displaystyle #1$}}%
  \else\hspace{2em}\phantom{#1}\hspace{2em}\fi
}

% 記述・計算スペース（薄い枠で高さを確保）
% 使い方: \writespace[高さ]{解答例}
\newcommand{\writespace}[2][5em]{%
  \begin{tcolorbox}[enhanced, colback=white, colframe=black!35,
      boxrule=0.4pt, sharp corners, height=#1, valign=center,
      left=8pt, right=8pt]
    \ifshowanswer\answerstyle{#2}\fi
  \end{tcolorbox}%
}

% ========================================
% 定義（赤系）・公式（青系）・定理（青系・太枠）
% 使い方: \begin{definitionbox}{用語} ... \end{definitionbox}
% ========================================
\tcbset{
  mathbox/.style={
    enhanced, sharp corners,
    left=10pt, right=10pt, top=8pt, bottom=8pt,
    fonttitle=\bfseries,
  }
}
\newtcolorbox{definitionbox}[1]{
  mathbox, colback=red!5, colframe=red!70!black, boxrule=0.8pt,
  code={\refstepcounter{definition}},
  title={定義 \thedefinition\quad #1}
}
\newtcolorbox{formulabox}[1]{
  mathbox, colback=blue!5, colframe=blue!70!black, boxrule=0.8pt,
  code={\refstepcounter{formula}},
  title={公式 \theformula\quad #1}
}
\newtcolorbox{theorembox}[1]{
  mathbox, colback=blue!5, colframe=blue!70!black, boxrule=1.2pt,
  code={\refstepcounter{theorem}},
  title={定理 \thetheorem\quad #1}
}
% テクニック（緑系）
% 使い方: \begin{techniquebox}[ラベル]{見出し} ... \end{techniquebox}
\newcounter{technique}[section]
\renewcommand{\thetechnique}{\arabic{mathchapter}.\arabic{section}.\arabic{technique}}
\newtcolorbox{techniquebox}[2][]{
  mathbox, colback=green!5, colframe=green!70!black, boxrule=0.8pt,
  code={\refstepcounter{technique}\if\relax\detokenize{#1}\relax\else\label{#1}\fi},
  title={Technique \thetechnique\if\relax\detokenize{#2}\relax\else\quad\textbf{#2}\fi}
}


% A4 2ページに収めるため，共通設定（1.25）より行間を少し詰める
\linespread{1.15}
\pagestyle{empty}

% ========================================
% ヘッダ（1ページ目の先頭）
% 使い方: \lectureheader{科目}{章番号}{章の名前}{節番号}{節の名前}{項番号-小単元}{小単元の名前}
%           {教科書のページ}{本時の目標}
%   例: \lectureheader{数学Ⅱ}{5}{指数関数と対数関数}{1}{指数関数}{2-A}{指数関数 $y=a^x$ のグラフ}
%         {150--151}{指数関数のグラフを描けるようになろう}
%   項番号-小単元は「2-A」のほか「2」だけでもよい（カウンターには先頭の数字を使う）
%   教科書のページを空 {} にすると，手書き用の空欄になる
%   タイトルと教科書のページが1行に収まらないときは，タイトルを横に縮めて収める
% ========================================
\newsavebox{\lectitlebox}
\newsavebox{\lecpagebox}
\newlength{\lectitlewidth}
\def\lecitemnumber#1-#2\relax{#1}
\newcommand{\lectureheader}[9]{%
  \setcounter{mathchapter}{#2}%
  \setcounter{section}{#4}%
  \setcounter{subsection}{\lecitemnumber#6-\relax}%
  \noindent{\bfseries #1}\quad 第#2章\quad #3\hfill
  \rule{0pt}{1.4em}%
  \underline{\hspace{1.5zw}}年\underline{\hspace{1.5zw}}組\underline{\hspace{2zw}}番\quad
  氏名\underline{\hspace{8zw}}\par
  \vspace{0.3em}
  \sbox{\lectitlebox}{\boldmath{\Large\bfseries 第#4節\enspace #5}\hspace{1zw}{\large\bfseries #6\enspace #7}}%
  \sbox{\lecpagebox}{教科書 p.\,\ifx\relax#8\relax\hspace*{2zw}\else#8\fi}%
  \setlength{\lectitlewidth}{\dimexpr\linewidth-\wd\lecpagebox-1zw\relax}%
  \noindent
  \ifdim\wd\lectitlebox>\lectitlewidth
    \resizebox{\lectitlewidth}{!}{\usebox{\lectitlebox}}%
  \else
    \usebox{\lectitlebox}%
  \fi
  \hfill\usebox{\lecpagebox}\par
  \vspace{-0.3em}
  \noindent\rule{\linewidth}{0.8pt}\par
  \vspace{0.3em}
  \begin{itembox}[l]{\textbf{本時の目標}}
    \centering\bfseries\boldmath #9
  \end{itembox}
  \vspace{0.3em}
}

% 次のページへ移る（1→2 は表面の右，2→3 は裏面の左，3→4 は裏面の右）
\newcommand{\nextpage}{\newpage}

% 小単元（教科書の「A」「B」…）: \topic{A}{平均変化率}
\newcommand{\topic}[2]{%
  \par\vspace{0.3em}\noindent{\Large\bfseries #1\quad #2}\par\vspace{0.3em}%
}

% 見出し: \heading{見出し}
\newcommand{\heading}[1]{%
  \par\noindent{\large\bfseries #1}\par\vspace{0.3em}%
}

% 前時の復習（導入）: \begin{reviewbox}{中学数学：1次関数の変化の割合} ... \end{reviewbox}
\newenvironment{reviewbox}[1]{%
  \begin{itembox}[l]{\textbf{前時の復習}\quad #1}
}{\end{itembox}}

% 本時のまとめ
\newenvironment{summarybox}{%
  \begin{itembox}[l]{\textbf{本時のまとめ}}
}{\end{itembox}}

% 左右配置（左：説明・計算，右：図）
% 使い方: \sidebyside[左の幅の割合]{左}{右}
\newcommand{\sidebyside}[3][0.5]{%
  \par\noindent
  \begin{minipage}[t]{#1\linewidth}\vspace{0pt}#2\end{minipage}\hfill
  \begin{minipage}[t]{\dimexpr0.98\linewidth-#1\linewidth\relax}%
    \vspace{0pt}\centering #3\end{minipage}\par
}

% 次回予告（最終ページの下端）
\newcommand{\nexttime}[1]{%
  \par\noindent\parbox{\linewidth}{%
    \rule{\linewidth}{0.4pt}\par
    次回：#1}\par
}

% ※ B4 に縮小（約87%）すると，14pt は実寸で約12pt になる
% ※ jsarticle は \mag で拡大するので，紙面の寸法は truemm で指定する．
%    TikZ の cm も拡大される（14pt では 1cm → 実寸 1.4cm）点に注意
% ========================================
% a4paper は \mag で拡大されてしまうため，用紙も truemm で実寸指定する
\usepackage[dvipdfm, paperwidth=210truemm, paperheight=297truemm,
            margin=15truemm]{geometry}

% ========================================
% 共通プリアンブル（講義MODE・演習MODE）
% パッケージ／解答切り替え／番号／空欄／定義・公式・定理の囲み
% ========================================
\usepackage{amsmath, amssymb, mathtools, bm}
\usepackage{ascmac}
\usepackage[dvipdfmx]{graphicx}
\usepackage[dvipdfmx]{xcolor}
\usepackage{tikz}
\usetikzlibrary{calc, arrows.meta}
\usepackage[most]{tcolorbox}
\usepackage{enumitem}
\usepackage{tasks}

\linespread{1.25}
\everymath{\displaystyle}
\setlength{\parindent}{0pt}
\setlist[enumerate,1]{label=（\arabic*）, leftmargin=*, labelsep=0.5em}

% ========================================
% 解答表示の切り替え（既定は教師用．make student で \studentmode が定義され生徒用になる）
% ========================================
\newif\ifshowanswer
\showanswertrue
\ifdefined\studentmode\showanswerfalse\fi

% ========================================
% 番号（「章.節.番号」形式，節ごとにリセット）
% ========================================
% jsarticle には章がないので，章番号は専用カウンタで持つ
\newcounter{mathchapter}
\newcounter{definition}[section]
\newcounter{formula}[section]
\newcounter{theorem}[section]
\renewcommand{\thedefinition}{\arabic{mathchapter}.\arabic{section}.\arabic{definition}}
\renewcommand{\theformula}{\arabic{mathchapter}.\arabic{section}.\arabic{formula}}
\renewcommand{\thetheorem}{\arabic{mathchapter}.\arabic{section}.\arabic{theorem}}

% ========================================
% 書き込み用コマンド
% ========================================
% 解答は青字・太字（数式も和文も太字になるよう \bfseries\boldmath を使う）
% 色は外側から \textcolor で掛けると数式中の \text{} の和文に効かないため，
% 箱の内側で \color を指定する
\newcommand{\answerstyle}[1]{{\color{blue}\bfseries\boldmath #1}}

% 空欄の左右それぞれの余白
\newcommand{\fitblankpad}{2zw}

% 空欄（答えの幅＋左右の余白で下線を引く）
% 使い方: \fitblankbf[]{答え}  ※文字列は \text{...} で囲む
\newsavebox{\fitblanksavebox}
\newcommand{\fitblankbf}[2][]{%
  \sbox{\fitblanksavebox}{\answerstyle{$\displaystyle #2$}}%
  \underline{\makebox[\dimexpr\wd\fitblanksavebox+\fitblankpad*2\relax][c]{%
    \ifshowanswer\usebox{\fitblanksavebox}%
    \else\phantom{\usebox{\fitblanksavebox}}\fi}}%
}

% 下線なしの空欄（分数の分子・分母など，下線と分数線が重なる場所で使う）
% 使い方: \frac{\fitblankbox{y \text{ の変化量}}}{\fitblankbox{x \text{ の変化量}}}
\newcommand{\fitblankbox}[1]{%
  \sbox{\fitblanksavebox}{\answerstyle{$\displaystyle #1$}}%
  \makebox[\dimexpr\wd\fitblanksavebox+\fitblankpad*2\relax][c]{%
    \ifshowanswer\usebox{\fitblanksavebox}%
    \else\phantom{\usebox{\fitblanksavebox}}\fi}%
}

% 数式中の計算欄（下線なし，答えの幅＋左右2em を確保）
% 使い方: $= \answermath{3 - 1 = 2}$
\newcommand{\answermath}[1]{%
  \ifshowanswer\mbox{\answerstyle{$\displaystyle #1$}}%
  \else\hspace{2em}\phantom{#1}\hspace{2em}\fi
}

% 記述・計算スペース（薄い枠で高さを確保）
% 使い方: \writespace[高さ]{解答例}
\newcommand{\writespace}[2][5em]{%
  \begin{tcolorbox}[enhanced, colback=white, colframe=black!35,
      boxrule=0.4pt, sharp corners, height=#1, valign=center,
      left=8pt, right=8pt]
    \ifshowanswer\answerstyle{#2}\fi
  \end{tcolorbox}%
}

% ========================================
% 定義（赤系）・公式（青系）・定理（青系・太枠）
% 使い方: \begin{definitionbox}{用語} ... \end{definitionbox}
% ========================================
\tcbset{
  mathbox/.style={
    enhanced, sharp corners,
    left=10pt, right=10pt, top=8pt, bottom=8pt,
    fonttitle=\bfseries,
  }
}
\newtcolorbox{definitionbox}[1]{
  mathbox, colback=red!5, colframe=red!70!black, boxrule=0.8pt,
  code={\refstepcounter{definition}},
  title={定義 \thedefinition\quad #1}
}
\newtcolorbox{formulabox}[1]{
  mathbox, colback=blue!5, colframe=blue!70!black, boxrule=0.8pt,
  code={\refstepcounter{formula}},
  title={公式 \theformula\quad #1}
}
\newtcolorbox{theorembox}[1]{
  mathbox, colback=blue!5, colframe=blue!70!black, boxrule=1.2pt,
  code={\refstepcounter{theorem}},
  title={定理 \thetheorem\quad #1}
}
% テクニック（緑系）
% 使い方: \begin{techniquebox}[ラベル]{見出し} ... \end{techniquebox}
\newcounter{technique}[section]
\renewcommand{\thetechnique}{\arabic{mathchapter}.\arabic{section}.\arabic{technique}}
\newtcolorbox{techniquebox}[2][]{
  mathbox, colback=green!5, colframe=green!70!black, boxrule=0.8pt,
  code={\refstepcounter{technique}\if\relax\detokenize{#1}\relax\else\label{#1}\fi},
  title={Technique \thetechnique\if\relax\detokenize{#2}\relax\else\quad\textbf{#2}\fi}
}


% A4 2ページに収めるため，共通設定（1.25）より行間を少し詰める
\linespread{1.15}
\pagestyle{empty}

% ========================================
% ヘッダ（1ページ目の先頭）
% 使い方: \lectureheader{科目}{章番号}{章の名前}{節番号}{節の名前}{項番号-小単元}{小単元の名前}
%           {教科書のページ}{本時の目標}
%   例: \lectureheader{数学Ⅱ}{5}{指数関数と対数関数}{1}{指数関数}{2-A}{指数関数 $y=a^x$ のグラフ}
%         {150--151}{指数関数のグラフを描けるようになろう}
%   項番号-小単元は「2-A」のほか「2」だけでもよい（カウンターには先頭の数字を使う）
%   教科書のページを空 {} にすると，手書き用の空欄になる
%   タイトルと教科書のページが1行に収まらないときは，タイトルを横に縮めて収める
% ========================================
\newsavebox{\lectitlebox}
\newsavebox{\lecpagebox}
\newlength{\lectitlewidth}
\def\lecitemnumber#1-#2\relax{#1}
\newcommand{\lectureheader}[9]{%
  \setcounter{mathchapter}{#2}%
  \setcounter{section}{#4}%
  \setcounter{subsection}{\lecitemnumber#6-\relax}%
  \noindent{\bfseries #1}\quad 第#2章\quad #3\hfill
  \rule{0pt}{1.4em}%
  \underline{\hspace{1.5zw}}年\underline{\hspace{1.5zw}}組\underline{\hspace{2zw}}番\quad
  氏名\underline{\hspace{8zw}}\par
  \vspace{0.3em}
  \sbox{\lectitlebox}{\boldmath{\Large\bfseries 第#4節\enspace #5}\hspace{1zw}{\large\bfseries #6\enspace #7}}%
  \sbox{\lecpagebox}{教科書 p.\,\ifx\relax#8\relax\hspace*{2zw}\else#8\fi}%
  \setlength{\lectitlewidth}{\dimexpr\linewidth-\wd\lecpagebox-1zw\relax}%
  \noindent
  \ifdim\wd\lectitlebox>\lectitlewidth
    \resizebox{\lectitlewidth}{!}{\usebox{\lectitlebox}}%
  \else
    \usebox{\lectitlebox}%
  \fi
  \hfill\usebox{\lecpagebox}\par
  \vspace{-0.3em}
  \noindent\rule{\linewidth}{0.8pt}\par
  \vspace{0.3em}
  \begin{itembox}[l]{\textbf{本時の目標}}
    \centering\bfseries\boldmath #9
  \end{itembox}
  \vspace{0.3em}
}

% 次のページへ移る（1→2 は表面の右，2→3 は裏面の左，3→4 は裏面の右）
\newcommand{\nextpage}{\newpage}

% 小単元（教科書の「A」「B」…）: \topic{A}{平均変化率}
\newcommand{\topic}[2]{%
  \par\vspace{0.3em}\noindent{\Large\bfseries #1\quad #2}\par\vspace{0.3em}%
}

% 見出し: \heading{見出し}
\newcommand{\heading}[1]{%
  \par\noindent{\large\bfseries #1}\par\vspace{0.3em}%
}

% 前時の復習（導入）: \begin{reviewbox}{中学数学：1次関数の変化の割合} ... \end{reviewbox}
\newenvironment{reviewbox}[1]{%
  \begin{itembox}[l]{\textbf{前時の復習}\quad #1}
}{\end{itembox}}

% 本時のまとめ
\newenvironment{summarybox}{%
  \begin{itembox}[l]{\textbf{本時のまとめ}}
}{\end{itembox}}

% 左右配置（左：説明・計算，右：図）
% 使い方: \sidebyside[左の幅の割合]{左}{右}
\newcommand{\sidebyside}[3][0.5]{%
  \par\noindent
  \begin{minipage}[t]{#1\linewidth}\vspace{0pt}#2\end{minipage}\hfill
  \begin{minipage}[t]{\dimexpr0.98\linewidth-#1\linewidth\relax}%
    \vspace{0pt}\centering #3\end{minipage}\par
}

% 次回予告（最終ページの下端）
\newcommand{\nexttime}[1]{%
  \par\noindent\parbox{\linewidth}{%
    \rule{\linewidth}{0.4pt}\par
    次回：#1}\par
}
